% Supplementary Note 4: the image check. Numbers through macros only.
\suppnote{The image check}

This note describes the check of the stimulus images that the main text's Results summarize in one sentence, in the section on how ratings of second works follow the level of modification.

Studies 2, 4 and 5 showed an original beside three second works that Alex was said to have made from it: None, the original shown again unchanged; Slight, a light edit; and Dramatic, a heavy one. Participants rated each second work's transformativeness, essence change and creativity. The design takes three things about the pictures for granted: that the ten originals differ from one another, that each artwork's Slight version is closer to it than its Dramatic version is, and that every edited version, however much it was changed, stays closer to its own original than to any of the other nine. The third matters most, because a question about how far a second work moved from its original rests on a pairing that the pictures must support. No registration names this check. It was specified on 10 October 2026, after every registered result was known, so it is exploratory and changes no registered verdict; its comparison with the ratings is descriptive, with no test.

\paragraph*{What was compared.}
The check used the image files the surveys displayed: the ten originals and their ten Slight and ten Dramatic versions. They were taken from a printout of the Study~5 survey, a PDF file printed from the survey software that holds each image as the survey served it, and the survey files of Studies 2, 4 and 5 name the same image files. The None second work is the original's own file, with identical pixels, so its distance from the original is zero. Distances came from three measures that notice different things. Two are image-recognition networks, programs trained on large collections of pictures that turn any image into a long list of numbers, so that any two images lie at a distance from each other.\cite{oquab2024dinov2,radford2021learning} DINOv2 learned, without labels, to give matching numbers to different views of the same photograph, so its numbers describe layout, shapes and content. CLIP learned to match pictures with their captions, so its numbers reflect what a picture would be described as. For both, the distance between two images is 1 minus the cosine similarity of their lists, a standard measure of how closely two lists of numbers point the same way, so identical images lie at zero. The third measure uses no network. It compares the mix of colors in two images, taken as the share of pixels that fall into each cell of a fixed grid of colors, and is zero when two images have the same mix. Each network saw each image in two ways: the whole picture resized to a square, the primary input, and the network's own standard crop, which keeps the central square and drops the edges of a picture that is not square. The rule for passing each check was written into the code before any distance was computed, and no person or language model judged any image.

\paragraph*{The three checks.}
On both networks and with both inputs, all three checks passed. The two closest originals lay farther apart than any original lay from its own Slight version. The Slight version was closer to its original than the Dramatic version for every artwork (\XSImageCheckTwoDINO{} artworks on DINOv2 and \XSImageCheckTwoCLIP{} on CLIP). These steps are ordered within each artwork but overlap across artworks: the Dramatic version of one artwork can lie nearer its original than the Slight version of another lies to its own. Every edited image lay closest to its own original (\XSImageCheckThreeDINO{} images on DINOv2 and \XSImageCheckThreeCLIP{} on CLIP). The two inputs gave the same verdict on every check.

The color measure failed all three checks. Some Slight edits moved an image's colors farther than the two most similar originals lie apart, so the first check failed. It put the Slight version closer than the Dramatic version for \XSImageCheckTwoColour{} artworks, not all ten, and put only half of the twenty edited images (\XSImageCheckThreeColour) closest to their own original. Because this measure ignores where colors lie in a picture, its failures show that the Slight edits changed palettes substantially and that some edited images have a palette nearer another artwork's. They do not show that an edited image resembles another artwork. The color measure counts a Slight edit as a much larger share of a Dramatic edit than either network does.

\paragraph*{The Slight versions were also cropped or resized.}
The registrations describe the Slight edits as adjustments of color and contrast, and the CogSci 2025 paper, our proceedings paper on Studies 2 and 4, describes them as adjustments of hue, coloration and contrast.\cite{kamper2025originalsin} But \XSImageCheckSlightReframedN{} of the ten Slight versions also differ from their original in pixel size and in the ratio of width to height, so those edits cropped or resized the picture as well (main-text Fig.~1a shows a Slight version that was not). The descriptions participants read never said how much had been changed. Supplementary Note~6 lists this cropping among the possible reasons for the peak in the label's penalty at Slight.

\paragraph*{Distance and ratings, description only.}
The check also set each edited image's distance from its original beside its mean rating in Studies 2, 4 and 5. Each image's mean was computed within each study and then averaged over the three, so each study counts once whatever its size. Across the twenty edited images (nineteen for transformativeness, because one Slight item used a six-point scale and was set aside), images farther from their original were rated higher, strongly so for every distance measure and every rating, by Spearman correlations, ordinary correlations computed on the ranks of the values. Most of that relation comes from the level of modification itself: every Dramatic image was rated above every Slight image, and the Dramatic images are as a group also farther from their originals. Within a level, across its ten images, the correlations were mostly small and changed with the measure, the input and the label group. With so few images these correlations describe the stimuli and test nothing, and no $p$-value is reported. The ratings are therefore shown to follow the level of modification; they are not shown to track the amount of visible change within a level.

\paragraph*{Limits.}
Neither network was trained on judgments of art or of artistic change, so neither distance measures what participants saw or judged; the checks ask whether the pictures are arranged as the design assumes, as two networks with different training describe them. The pixels are those of the printout. The survey files show that Studies 2, 4 and 5 named the same images, but they cannot show that the survey platform served the same bytes on every date each study ran, and no copy of the live survey pages was kept. That the None second work showed the original rests on the survey file, the printout and the ratings near the floor of the scale at None.

\paragraph*{Code and images.}
The image check's folder in the analysis code (main text, Code availability) holds the script that extracts the images, the analysis notebook, the images themselves with a checksum for each (a fingerprint of a file's exact bytes), the pinned versions of both networks, every table and two figures. Its README explains how to rerun the check from a fresh copy. The networks run on one processor thread with a deterministic setting, so a rerun on the same kind of computer gives the same numbers.
